\documentclass[../main.tex]{subfiles}
\begin{document}

\chapter{Synchronization Constructs}
\lessoninfo{
  video={P3L4},
  textbook={OSTEP \cite{ostep} ch.~28, 31; Silberschatz et al.\ \cite{silberschatz2018} ch.~6--7},
  keywords={spinlock, busy waiting, semaphore, reader-writer lock, monitor, serializer, path expression, barrier, wait-free synchronization, RCU, atomic instruction, test-and-set, cache coherence, test-and-test-and-set},
}

Lesson~3 introduced mutexes and condition variables, and Lesson~9 showed
that processes communicating through shared memory need synchronization
just as threads do.  This lesson surveys further synchronization
constructs---spinlocks, semaphores, reader-writer locks, monitors and
several less common ones---that are either simpler or more expressive than
mutexes and condition variables.  It then turns to what every such
construct ultimately rests on: atomic instructions provided by the
hardware, and their interaction with the caches of a shared memory
multiprocessor.

\section{Beyond mutexes and condition variables}
\label{sec:l10-limits}

Mutexes and condition variables can express the synchronization policies
of Lesson~3, but they have three limitations that motivate other
constructs.\source{P3L4, more about synchronization; OSTEP ch.~28.}
\begin{description}
  \item[Error-prone use] Correctness rests entirely on the programmer.
    Every access to shared state must lock the right mutex and release it
    on every path, and every change of state must signal or broadcast the
    right condition variable.  Most pitfalls listed in Lesson~3, such as
    unlocking the wrong mutex or signalling the wrong condition variable,
    violate these rules, and the constructs themselves detect none of
    them.
  \item[Limited expressive power] A mutex enforces exactly one policy: one
    thread at a time.  A richer policy, such as admitting several readers
    or one writer, or giving waiting writers precedence, must be encoded
    with helper variables and explicit checks, like the proxy variable of
    the readers-writers solution in Lesson~3.
  \item[Dependence on low-level support] A mutex is itself shared state.
    Its lock operation must find the mutex free and mark it as owned in one
    indivisible step, which ordinary instructions cannot guarantee; it
    needs atomic instructions from the hardware (\cref{sec:l10-hw}).
\end{description}
The first two limitations motivate the constructs of
\cref{sec:l10-spinlock,sec:l10-semaphore,sec:l10-rwlock,sec:l10-monitor,sec:l10-other};
the third is the subject of the second half of the lesson.

\section{Spinlocks}
\label{sec:l10-spinlock}

The simplest construct beyond the mutex is the spinlock.\source{P3L4,
spinlocks; OSTEP ch.~28.}

\begin{definition}[Spinlock]
  A \term{spinlock}[スピンロック] is a lock that provides mutual exclusion
  through lock and unlock operations, like a mutex, except that a thread
  which finds the lock held does not block.  The thread keeps running and
  repeatedly checks the lock until the lock becomes free or the thread is
  preempted.
\end{definition}

Checking a condition over and over in a loop, instead of giving up the CPU
until the condition holds, is called \term{busy waiting}[ビジーウェイト]; a
thread that busy-waits on a spinlock is said to \emph{spin}.  A spinlock is
used exactly like a mutex:
\begin{lstlisting}[language=C, numbers=none]
spinlock_lock(s);
// critical section
spinlock_unlock(s);
\end{lstlisting}

The two constructs differ only in what a thread does while the lock is held
by another thread, and they incur different costs
(\cref{fig:l10-spin-block}).  A thread that blocks on a mutex consumes no
CPU time while it waits, but it pays for a context switch away from it and,
once the mutex is released, for a wake-up and a context switch back.  A
spinning thread avoids both context switches, but it occupies its CPU for
as long as the lock remains held.\margin{POSIX threads provide spinlocks
as \texttt{pthread\_spinlock\_t} with \texttt{pthread\_spin\_lock} and
\texttt{pthread\_spin\_unlock}.}

\begin{figure}[htbp]
  \centering
  % Blocking on a mutex versus spinning on a spinlock: timelines of two CPUs.
% Usage: % Blocking on a mutex versus spinning on a spinlock: timelines of two CPUs.
% Usage: % Blocking on a mutex versus spinning on a spinlock: timelines of two CPUs.
% Usage: \input{figures/l10-spin-block}   (width ~116 mm)
\begin{tikzpicture}[x=1mm, y=1mm, font=\footnotesize\sffamily,
  lab/.style={font=\scriptsize\sffamily, align=center},
  row/.style={font=\scriptsize\sffamily, anchor=east},
  ttl/.style={font=\footnotesize\sffamily\bfseries, anchor=west},
  >={Stealth[length=1.6mm]}]
  % \ltenseg{x1}{x2}{y}{fill}{text}
  \def\ltenseg#1#2#3#4#5{%
    \draw[fill=#4] (#1,#3-2.5) rectangle (#2,#3+2.5);
    \node[lab] at ({(#1+#2)/2},#3) {#5};}
  % event lines
  \draw[densely dashed, ln-muted] (50,33) -- (50,-12);
  \draw[densely dashed, ln-muted] (76,33) -- (76,-12);
  \node[lab, anchor=south, fill=white, inner sep=1pt] at (50,33)
    {\iflangja{T2 が lock を呼ぶ}{T2 calls lock}};
  \node[lab, anchor=south, fill=white, inner sep=1pt] at (76,33)
    {\iflangja{T1 が unlock する}{T1 unlocks}};
  % (a) mutex
  \node[ttl] at (0,29) {\iflangja{(a) ミューテックス：T2 はブロック}{(a) mutex: T2 blocks}};
  \node[row] at (22,22) {CPU 1};
  \ltenseg{24}{36}{22}{white}{T1}
  \ltenseg{36}{76}{22}{ln-box}{\iflangja{T1：クリティカルセクション}{T1: critical section}}
  \ltenseg{76}{116}{22}{white}{T1}
  \node[row] at (22,15) {CPU 2};
  \ltenseg{24}{50}{15}{white}{T2}
  \ltenseg{50}{55}{15}{black!35}{}
  \ltenseg{55}{76}{15}{white}{\iflangja{T3（別のスレッド）}{T3 (another thread)}}
  \ltenseg{76}{81}{15}{black!35}{}
  \ltenseg{81}{116}{15}{ln-box}{\iflangja{T2：クリティカルセクション}{T2: critical section}}
  % (b) spinlock
  \node[ttl] at (0,6) {\iflangja{(b) スピンロック：T2 はスピン}{(b) spinlock: T2 spins}};
  \node[row] at (22,-1) {CPU 1};
  \ltenseg{24}{36}{-1}{white}{T1}
  \ltenseg{36}{76}{-1}{ln-box}{\iflangja{T1：クリティカルセクション}{T1: critical section}}
  \ltenseg{76}{116}{-1}{white}{T1}
  \node[row] at (22,-8) {CPU 2};
  \ltenseg{24}{50}{-8}{white}{T2}
  \ltenseg{50}{76}{-8}{ln-caution!20}{\iflangja{ビジーウェイト}{busy waiting}}
  \ltenseg{76}{116}{-8}{ln-box}{\iflangja{T2：クリティカルセクション}{T2: critical section}}
  % legend
  \draw[fill=black!35] (24,-18) rectangle (28,-15);
  \node[lab, anchor=west] at (29,-16.5) {\iflangja{コンテキストスイッチ}{context switch}};
  \draw[fill=ln-box] (64,-18) rectangle (68,-15);
  \node[lab, anchor=west] at (69,-16.5) {\iflangja{ロックを保持}{lock held}};
  \draw[->] (92,-16.5) -- node[lab, above] {\iflangja{時間}{time}} (116,-16.5);
\end{tikzpicture}
   (width ~116 mm)
\begin{tikzpicture}[x=1mm, y=1mm, font=\footnotesize\sffamily,
  lab/.style={font=\scriptsize\sffamily, align=center},
  row/.style={font=\scriptsize\sffamily, anchor=east},
  ttl/.style={font=\footnotesize\sffamily\bfseries, anchor=west},
  >={Stealth[length=1.6mm]}]
  % \ltenseg{x1}{x2}{y}{fill}{text}
  \def\ltenseg#1#2#3#4#5{%
    \draw[fill=#4] (#1,#3-2.5) rectangle (#2,#3+2.5);
    \node[lab] at ({(#1+#2)/2},#3) {#5};}
  % event lines
  \draw[densely dashed, ln-muted] (50,33) -- (50,-12);
  \draw[densely dashed, ln-muted] (76,33) -- (76,-12);
  \node[lab, anchor=south, fill=white, inner sep=1pt] at (50,33)
    {\iflangja{T2 が lock を呼ぶ}{T2 calls lock}};
  \node[lab, anchor=south, fill=white, inner sep=1pt] at (76,33)
    {\iflangja{T1 が unlock する}{T1 unlocks}};
  % (a) mutex
  \node[ttl] at (0,29) {\iflangja{(a) ミューテックス：T2 はブロック}{(a) mutex: T2 blocks}};
  \node[row] at (22,22) {CPU 1};
  \ltenseg{24}{36}{22}{white}{T1}
  \ltenseg{36}{76}{22}{ln-box}{\iflangja{T1：クリティカルセクション}{T1: critical section}}
  \ltenseg{76}{116}{22}{white}{T1}
  \node[row] at (22,15) {CPU 2};
  \ltenseg{24}{50}{15}{white}{T2}
  \ltenseg{50}{55}{15}{black!35}{}
  \ltenseg{55}{76}{15}{white}{\iflangja{T3（別のスレッド）}{T3 (another thread)}}
  \ltenseg{76}{81}{15}{black!35}{}
  \ltenseg{81}{116}{15}{ln-box}{\iflangja{T2：クリティカルセクション}{T2: critical section}}
  % (b) spinlock
  \node[ttl] at (0,6) {\iflangja{(b) スピンロック：T2 はスピン}{(b) spinlock: T2 spins}};
  \node[row] at (22,-1) {CPU 1};
  \ltenseg{24}{36}{-1}{white}{T1}
  \ltenseg{36}{76}{-1}{ln-box}{\iflangja{T1：クリティカルセクション}{T1: critical section}}
  \ltenseg{76}{116}{-1}{white}{T1}
  \node[row] at (22,-8) {CPU 2};
  \ltenseg{24}{50}{-8}{white}{T2}
  \ltenseg{50}{76}{-8}{ln-caution!20}{\iflangja{ビジーウェイト}{busy waiting}}
  \ltenseg{76}{116}{-8}{ln-box}{\iflangja{T2：クリティカルセクション}{T2: critical section}}
  % legend
  \draw[fill=black!35] (24,-18) rectangle (28,-15);
  \node[lab, anchor=west] at (29,-16.5) {\iflangja{コンテキストスイッチ}{context switch}};
  \draw[fill=ln-box] (64,-18) rectangle (68,-15);
  \node[lab, anchor=west] at (69,-16.5) {\iflangja{ロックを保持}{lock held}};
  \draw[->] (92,-16.5) -- node[lab, above] {\iflangja{時間}{time}} (116,-16.5);
\end{tikzpicture}
   (width ~116 mm)
\begin{tikzpicture}[x=1mm, y=1mm, font=\footnotesize\sffamily,
  lab/.style={font=\scriptsize\sffamily, align=center},
  row/.style={font=\scriptsize\sffamily, anchor=east},
  ttl/.style={font=\footnotesize\sffamily\bfseries, anchor=west},
  >={Stealth[length=1.6mm]}]
  % \ltenseg{x1}{x2}{y}{fill}{text}
  \def\ltenseg#1#2#3#4#5{%
    \draw[fill=#4] (#1,#3-2.5) rectangle (#2,#3+2.5);
    \node[lab] at ({(#1+#2)/2},#3) {#5};}
  % event lines
  \draw[densely dashed, ln-muted] (50,33) -- (50,-12);
  \draw[densely dashed, ln-muted] (76,33) -- (76,-12);
  \node[lab, anchor=south, fill=white, inner sep=1pt] at (50,33)
    {\iflangja{T2 が lock を呼ぶ}{T2 calls lock}};
  \node[lab, anchor=south, fill=white, inner sep=1pt] at (76,33)
    {\iflangja{T1 が unlock する}{T1 unlocks}};
  % (a) mutex
  \node[ttl] at (0,29) {\iflangja{(a) ミューテックス：T2 はブロック}{(a) mutex: T2 blocks}};
  \node[row] at (22,22) {CPU 1};
  \ltenseg{24}{36}{22}{white}{T1}
  \ltenseg{36}{76}{22}{ln-box}{\iflangja{T1：クリティカルセクション}{T1: critical section}}
  \ltenseg{76}{116}{22}{white}{T1}
  \node[row] at (22,15) {CPU 2};
  \ltenseg{24}{50}{15}{white}{T2}
  \ltenseg{50}{55}{15}{black!35}{}
  \ltenseg{55}{76}{15}{white}{\iflangja{T3（別のスレッド）}{T3 (another thread)}}
  \ltenseg{76}{81}{15}{black!35}{}
  \ltenseg{81}{116}{15}{ln-box}{\iflangja{T2：クリティカルセクション}{T2: critical section}}
  % (b) spinlock
  \node[ttl] at (0,6) {\iflangja{(b) スピンロック：T2 はスピン}{(b) spinlock: T2 spins}};
  \node[row] at (22,-1) {CPU 1};
  \ltenseg{24}{36}{-1}{white}{T1}
  \ltenseg{36}{76}{-1}{ln-box}{\iflangja{T1：クリティカルセクション}{T1: critical section}}
  \ltenseg{76}{116}{-1}{white}{T1}
  \node[row] at (22,-8) {CPU 2};
  \ltenseg{24}{50}{-8}{white}{T2}
  \ltenseg{50}{76}{-8}{ln-caution!20}{\iflangja{ビジーウェイト}{busy waiting}}
  \ltenseg{76}{116}{-8}{ln-box}{\iflangja{T2：クリティカルセクション}{T2: critical section}}
  % legend
  \draw[fill=black!35] (24,-18) rectangle (28,-15);
  \node[lab, anchor=west] at (29,-16.5) {\iflangja{コンテキストスイッチ}{context switch}};
  \draw[fill=ln-box] (64,-18) rectangle (68,-15);
  \node[lab, anchor=west] at (69,-16.5) {\iflangja{ロックを保持}{lock held}};
  \draw[->] (92,-16.5) -- node[lab, above] {\iflangja{時間}{time}} (116,-16.5);
\end{tikzpicture}

  \caption{Thread T2 requests a lock held by T1 on another CPU.  (a) On a
    mutex, T2 blocks, the CPU runs another thread, and T2 is switched back
    in after the unlock.  (b) On a spinlock, T2 busy-waits and enters the
    critical section as soon as T1 unlocks.}
  \label{fig:l10-spin-block}
\end{figure}

Spinning therefore pays off only when the lock is held for less time than
it takes to block and resume a thread, and only when the owner is running
on another CPU at the same time.  On a single CPU the owner cannot release
the lock while the spinning thread occupies the CPU, so spinning merely
wastes the remainder of the timeslice.  The adaptive mutexes of Lesson~5
combine both behaviours: they spin when the owner is running on another
CPU and block otherwise.\index{adaptive mutex}

\begin{example}
  Suppose blocking a thread and later resuming it costs \SI{8}{\micro\second}
  in total.  If the owner holds a lock for at most
  \SI{0.5}{\micro\second}, a waiting thread spins for at most
  \SI{0.5}{\micro\second}, one sixteenth of the cost of blocking.  If
  instead the owner performs a \SI{5}{\milli\second} disk read while
  holding the lock, a spinning thread burns \SI{5}{\milli\second} of CPU
  time, $5000/8 = 625$ times the cost of blocking.
\end{example}

Because a spinlock needs nothing but a shared word and an atomic
instruction, it is the basic construct from which more complex constructs
are built.  An operating system kernel, for example, protects the short
critical sections that update the state of a mutex and its queue of
blocked threads with a spinlock.

\section{Semaphores}
\label{sec:l10-semaphore}

Semaphores, proposed by Dijkstra \cite{dijkstra1965}, are a common
synchronization construct in operating system kernels.\source{P3L4,
semaphores, POSIX semaphores API; Dijkstra \cite{dijkstra1965}; OSTEP
ch.~31.}  They generalize the mutex from one thread at a time to a bounded
number of threads.

\begin{definition}[Semaphore]
  A \term{semaphore}[セマフォ] is a synchronization construct represented
  by an integer count, initialized with a non-negative integer, and two
  operations.  \emph{Wait}: if the count is positive, decrement it and
  proceed; if it is zero, block until it becomes positive.  \emph{Post}:
  increment the count and, if threads are blocked, wake one of them.
\end{definition}

The count records how many more threads may pass.  A thread decrements it
with wait when it enters and increments it with post when it leaves, so
the initial value is the maximum number of threads that can be inside at
the same time.  Dijkstra named the two operations P and V, after the Dutch
\emph{proberen} (to test) and \emph{verhogen} (to increase).

Semaphores are distinguished by the range of the count.  A
\term{counting semaphore}[計数セマフォ] may take any non-negative value and
admits as many threads at once as its initial value.  A
semaphore whose count takes only the values 0 and 1 is a
\term{binary semaphore}[バイナリセマフォ].  Initialized to 1, it admits
one thread at a time and behaves like a mutex.

\begin{example}
  A counting semaphore initialized to 2 guards two identical devices, and
  three threads use them.  \Cref{tab:l10-sem-trace} traces the count and
  the blocked threads.
\end{example}

\begin{table}[htbp]
  \centering
  \caption{A counting semaphore with initial value 2 used by three
    threads.  T3 blocks while the count is zero and proceeds after T1's
    post.}
  \label{tab:l10-sem-trace}
  \small
  \begin{tabular}{@{}clcl@{}}
    \toprule
    Step & Event & Count & Blocked \\
    \midrule
    1 & T1 waits: count $>0$, decrements, proceeds & 1 & --- \\
    2 & T2 waits: count $>0$, decrements, proceeds & 0 & --- \\
    3 & T3 waits: count $=0$, blocks               & 0 & T3 \\
    4 & T1 posts: increments, wakes T3             & 1 & --- \\
    5 & T3 re-tests: count $>0$, decrements, proceeds & 0 & --- \\
    6 & T2 posts                                    & 1 & --- \\
    7 & T3 posts                                    & 2 & --- \\
    \bottomrule
  \end{tabular}
\end{table}

\subsection{POSIX semaphores}

\needspace{9\baselineskip}
POSIX defines semaphores in \texttt{<semaphore.h>}:
\begin{lstlisting}[language=C, numbers=none]
#include <semaphore.h>

sem_t sem;
int sem_init(sem_t *sem, int pshared,
             unsigned int value);
int sem_wait(sem_t *sem);
int sem_post(sem_t *sem);
\end{lstlisting}
\texttt{sem\_init} sets the initial count to \texttt{value}.  The flag
\texttt{pshared} states whether the semaphore is used only by the threads of
one process (zero) or shared among processes (non-zero); in the latter case
the \texttt{sem\_t} itself must reside in memory that those processes share,
such as a shared memory segment of Lesson~9.  \texttt{sem\_wait} and
\texttt{sem\_post} are the wait and post operations.%
\caution{A semaphore has no owner: any thread may post, including one that
never waited.  A mutex is released by the thread that owns it.}

\needspace{14\baselineskip}
\begin{example}
  A server allows at most four requests at a time to reach its back-end
  database:
  \begin{lstlisting}[language=C, numbers=none]
#include <semaphore.h>

#define MAX_BACKEND 4
sem_t backend_slots;

void init_slots(void) {
  sem_init(&backend_slots, 0, MAX_BACKEND);
}

void handle_request(struct request *r) {
  sem_wait(&backend_slots);   /* take a slot */
  forward_to_backend(r);      /* at most 4 here */
  sem_post(&backend_slots);   /* free the slot */
}
  \end{lstlisting}
  A fifth thread that calls \texttt{handle\_request} while four requests
  are in progress blocks in \texttt{sem\_wait} until one of them posts.
\end{example}

\section{Reader-writer locks}
\label{sec:l10-rwlock}

Lesson~3 built the readers-writers policy---any number of readers or a
single writer---from a mutex, two condition variables and a proxy
variable.\source{P3L4, reader/writer locks, using reader/writer locks;
OSTEP ch.~31.}  Because the policy is so common, operating systems and
thread libraries offer it as a construct of its own.

\begin{definition}[Reader-writer lock]
  A \term{reader-writer lock}[リーダ・ライタロック] is a lock with two
  modes: \emph{read} (shared) mode, which any number of threads may hold at
  the same time, and \emph{write} (exclusive) mode, which one thread may
  hold only while no other thread holds the lock in either mode.
\end{definition}

\needspace{12\baselineskip}
A thread states in which mode it needs the lock, and the construct
enforces the policy.  In the Linux kernel the type is
\texttt{rwlock\_t}:\margin{Linux's \texttt{rwlock\_t} spins; the blocking
variant is \texttt{struct rw\_semaphore}.}
\begin{lstlisting}[language=C, numbers=none]
#include <linux/spinlock.h>

rwlock_t m;

read_lock(&m);
/* critical section: read shared state */
read_unlock(&m);

write_lock(&m);
/* critical section: read and modify */
write_unlock(&m);
\end{lstlisting}
POSIX threads, Java and the .NET framework of Windows provide the same
construct; in POSIX threads it reads:
\begin{lstlisting}[language=C, numbers=none]
pthread_rwlock_t m;
pthread_rwlock_rdlock(&m);   /* acquire in read mode */
pthread_rwlock_wrlock(&m);   /* acquire in write mode */
pthread_rwlock_unlock(&m);   /* release either mode */
\end{lstlisting}

Implementations agree on the basic policy but differ in details that affect
program behaviour:
\begin{itemize}
  \item \textbf{Recursive read locking.}  When a thread acquires the read
    lock twice, implementations differ in whether one
    \texttt{read\_unlock} releases both acquisitions or only one of them.
  \item \textbf{Upgrading and downgrading.}  Some implementations let a
    reader convert its lock into a write lock, or a writer convert its lock
    into a read lock, without releasing it; they differ in whether this is
    allowed and in the priority such a request has over threads already
    waiting.
  \item \textbf{Interaction with scheduling.}  An implementation may block
    an arriving reader while a writer, possibly of higher priority, is
    waiting, even though the lock is currently held only in read mode.
\end{itemize}

\begin{example}
  Readers arrive every \SI{3}{\milli\second} and each holds the lock in
  read mode for \SI{5}{\milli\second}.  Every reader arrives before its
  predecessor leaves, so the lock is never free.  An implementation that
  always admits a reader while the lock is held in read mode therefore
  starves a waiting writer indefinitely.  An implementation that blocks new
  readers once a writer waits admits the writer at most
  \SI{5}{\milli\second} after it starts waiting, as soon as the readers
  already inside have left.
\end{example}

\section{Monitors}
\label{sec:l10-monitor}

The constructs so far still leave it to the programmer to lock and unlock
at the right places.\source{P3L4, monitors; Hoare \cite{hoare1974};
Lampson and Redell \cite{lampson1980}; Silberschatz et al.\ ch.~6.}
Monitors move this responsibility into the construct.

\begin{definition}[Monitor]
  A \term{monitor}[モニタ] is a high-level synchronization construct that
  encapsulates a shared resource.  It specifies the resource, the
  procedures through which alone the resource may be accessed---each an
  \term{entry procedure}[エントリ手続き]---and the condition variables on
  which threads inside the monitor may wait.
\end{definition}

Mutual exclusion is implicit (\cref{fig:l10-monitor}).  When a thread calls
an entry procedure, the monitor's lock is acquired automatically, and
threads that find the monitor occupied wait on its entry queue; when the
procedure returns, the lock is released automatically.  A thread that waits
on a condition variable inside the monitor releases the lock, and when it
is woken it re-acquires the lock before it continues, again without any
explicit lock operation.  In the monitors of Hoare and of Mesa, the
conditions on which threads wait and the signal and broadcast operations
that wake them remain the programmer's responsibility; monitor designs with
automatic signalling go further and re-evaluate the waiting conditions
themselves whenever a thread leaves the monitor.

\begin{example}
  A bounded buffer of capacity \texttt{N} as a monitor, in C-like
  pseudocode:
  \begin{lstlisting}[language=C, numbers=none]
monitor BoundedBuffer {
  item buf[N];
  int count = 0, in = 0, out = 0;
  condition not_full, not_empty;

  entry void put(item x) {
    while (count == N)
      wait(not_full);
    buf[in] = x;  in = (in + 1) % N;  count++;
    signal(not_empty);
  }

  entry item get(void) {
    while (count == 0)
      wait(not_empty);
    item x = buf[out];  out = (out + 1) % N;
    count--;
    signal(not_full);
    return x;
  }
}
  \end{lstlisting}
  No lock operation appears in the code, so none can be forgotten or
  applied to the wrong lock.
\end{example}

\begin{figure}[htbp]
  \centering
  % Structure of a monitor: entry queue, implicit lock, entry procedures,
% shared data and condition-variable queues.
% Usage: % Structure of a monitor: entry queue, implicit lock, entry procedures,
% shared data and condition-variable queues.
% Usage: % Structure of a monitor: entry queue, implicit lock, entry procedures,
% shared data and condition-variable queues.
% Usage: \input{figures/l10-monitor}   (width ~116 mm)
\begin{tikzpicture}[x=1mm, y=1mm, font=\footnotesize\sffamily,
  box/.style={draw, rounded corners=2pt, fill=white, align=center, inner sep=2pt},
  q/.style={draw, rounded corners=2pt, fill=ln-box, minimum width=22mm,
            minimum height=7mm, inner sep=1pt},
  th/.style={draw, circle, fill=white, minimum size=4.2mm, inner sep=0pt,
             font=\tiny\sffamily},
  lab/.style={font=\scriptsize\sffamily, align=center},
  >={Stealth[length=1.8mm]}]
  % monitor
  \draw[thick, rounded corners=4pt, fill=ln-box] (32,-17) rectangle (84,19);
  \node[font=\footnotesize\sffamily\bfseries, anchor=north west] at (33,18.5)
    {\iflangja{モニタ}{monitor}};
  \node[box, minimum width=44mm, minimum height=7mm] (data) at (58,6)
    {\iflangja{共有データ}{shared data}\quad\texttt{buf}, \texttt{count}};
  \node[box, minimum width=19mm, minimum height=6mm, font=\footnotesize\ttfamily] (put) at (46.5,-4) {put()};
  \node[box, minimum width=19mm, minimum height=6mm, font=\footnotesize\ttfamily] (get) at (69.5,-4) {get()};
  \node[lab, anchor=north] at (58,-8.5) {\iflangja{エントリ手続き}{entry procedures}\\
    \iflangja{（内部で動くのは一度に1スレッド）}{(one thread inside at a time)}};
  \draw[ln-muted] (put.north) -- (data.south -| put.north);
  \draw[ln-muted] (get.north) -- (data.south -| get.north);
  % entry queue
  \node[q] (eq) at (13,0) {};
  \foreach \i/\t in {0/T5,1/T4} { \node[th] at (7.5+6*\i,0) {\t}; }
  \node[lab, anchor=south] at (13,3.5) {\iflangja{入口キュー}{entry queue}};
  \draw[->] (eq.east) -- node[lab, above, pos=0.45] {\iflangja{ロック}{lock}} (32,0);
  % condition queues
  \node[q] (nf) at (105,9) {};
  \node[th] at (99.5,9) {T2};
  \node[lab, anchor=south] at (105,12.5) {\texttt{not\_full}};
  \node[q] (ne) at (105,-9) {};
  \node[th] at (99.5,-9) {T3};
  \node[lab, anchor=south] at (105,-5.5) {\texttt{not\_empty}};
  \draw[->] (84,9) -- (nf.west);
  \draw[->] (84,-9) -- (ne.west);
  \node[lab, anchor=south] at (89,11.5) {\texttt{wait}};
  \node[lab] at (90,0) {\iflangja{ロック\\を解放}{releases\\lock}};
  % signal path back to the entry queue
  \draw[->] (ne.south) -- ++(0,-12) -| node[lab, pos=0.25, below]
    {\texttt{signal} / \texttt{broadcast}:
     \iflangja{起こされたスレッドはロックを再獲得する}{a woken thread re-acquires the lock}} (eq.south);
  % exit
  \draw[->] (58,19) -- ++(0,6) node[lab, anchor=south] {\iflangja{return（ロックを解放）}{return (releases lock)}};
\end{tikzpicture}
   (width ~116 mm)
\begin{tikzpicture}[x=1mm, y=1mm, font=\footnotesize\sffamily,
  box/.style={draw, rounded corners=2pt, fill=white, align=center, inner sep=2pt},
  q/.style={draw, rounded corners=2pt, fill=ln-box, minimum width=22mm,
            minimum height=7mm, inner sep=1pt},
  th/.style={draw, circle, fill=white, minimum size=4.2mm, inner sep=0pt,
             font=\tiny\sffamily},
  lab/.style={font=\scriptsize\sffamily, align=center},
  >={Stealth[length=1.8mm]}]
  % monitor
  \draw[thick, rounded corners=4pt, fill=ln-box] (32,-17) rectangle (84,19);
  \node[font=\footnotesize\sffamily\bfseries, anchor=north west] at (33,18.5)
    {\iflangja{モニタ}{monitor}};
  \node[box, minimum width=44mm, minimum height=7mm] (data) at (58,6)
    {\iflangja{共有データ}{shared data}\quad\texttt{buf}, \texttt{count}};
  \node[box, minimum width=19mm, minimum height=6mm, font=\footnotesize\ttfamily] (put) at (46.5,-4) {put()};
  \node[box, minimum width=19mm, minimum height=6mm, font=\footnotesize\ttfamily] (get) at (69.5,-4) {get()};
  \node[lab, anchor=north] at (58,-8.5) {\iflangja{エントリ手続き}{entry procedures}\\
    \iflangja{（内部で動くのは一度に1スレッド）}{(one thread inside at a time)}};
  \draw[ln-muted] (put.north) -- (data.south -| put.north);
  \draw[ln-muted] (get.north) -- (data.south -| get.north);
  % entry queue
  \node[q] (eq) at (13,0) {};
  \foreach \i/\t in {0/T5,1/T4} { \node[th] at (7.5+6*\i,0) {\t}; }
  \node[lab, anchor=south] at (13,3.5) {\iflangja{入口キュー}{entry queue}};
  \draw[->] (eq.east) -- node[lab, above, pos=0.45] {\iflangja{ロック}{lock}} (32,0);
  % condition queues
  \node[q] (nf) at (105,9) {};
  \node[th] at (99.5,9) {T2};
  \node[lab, anchor=south] at (105,12.5) {\texttt{not\_full}};
  \node[q] (ne) at (105,-9) {};
  \node[th] at (99.5,-9) {T3};
  \node[lab, anchor=south] at (105,-5.5) {\texttt{not\_empty}};
  \draw[->] (84,9) -- (nf.west);
  \draw[->] (84,-9) -- (ne.west);
  \node[lab, anchor=south] at (89,11.5) {\texttt{wait}};
  \node[lab] at (90,0) {\iflangja{ロック\\を解放}{releases\\lock}};
  % signal path back to the entry queue
  \draw[->] (ne.south) -- ++(0,-12) -| node[lab, pos=0.25, below]
    {\texttt{signal} / \texttt{broadcast}:
     \iflangja{起こされたスレッドはロックを再獲得する}{a woken thread re-acquires the lock}} (eq.south);
  % exit
  \draw[->] (58,19) -- ++(0,6) node[lab, anchor=south] {\iflangja{return（ロックを解放）}{return (releases lock)}};
\end{tikzpicture}
   (width ~116 mm)
\begin{tikzpicture}[x=1mm, y=1mm, font=\footnotesize\sffamily,
  box/.style={draw, rounded corners=2pt, fill=white, align=center, inner sep=2pt},
  q/.style={draw, rounded corners=2pt, fill=ln-box, minimum width=22mm,
            minimum height=7mm, inner sep=1pt},
  th/.style={draw, circle, fill=white, minimum size=4.2mm, inner sep=0pt,
             font=\tiny\sffamily},
  lab/.style={font=\scriptsize\sffamily, align=center},
  >={Stealth[length=1.8mm]}]
  % monitor
  \draw[thick, rounded corners=4pt, fill=ln-box] (32,-17) rectangle (84,19);
  \node[font=\footnotesize\sffamily\bfseries, anchor=north west] at (33,18.5)
    {\iflangja{モニタ}{monitor}};
  \node[box, minimum width=44mm, minimum height=7mm] (data) at (58,6)
    {\iflangja{共有データ}{shared data}\quad\texttt{buf}, \texttt{count}};
  \node[box, minimum width=19mm, minimum height=6mm, font=\footnotesize\ttfamily] (put) at (46.5,-4) {put()};
  \node[box, minimum width=19mm, minimum height=6mm, font=\footnotesize\ttfamily] (get) at (69.5,-4) {get()};
  \node[lab, anchor=north] at (58,-8.5) {\iflangja{エントリ手続き}{entry procedures}\\
    \iflangja{（内部で動くのは一度に1スレッド）}{(one thread inside at a time)}};
  \draw[ln-muted] (put.north) -- (data.south -| put.north);
  \draw[ln-muted] (get.north) -- (data.south -| get.north);
  % entry queue
  \node[q] (eq) at (13,0) {};
  \foreach \i/\t in {0/T5,1/T4} { \node[th] at (7.5+6*\i,0) {\t}; }
  \node[lab, anchor=south] at (13,3.5) {\iflangja{入口キュー}{entry queue}};
  \draw[->] (eq.east) -- node[lab, above, pos=0.45] {\iflangja{ロック}{lock}} (32,0);
  % condition queues
  \node[q] (nf) at (105,9) {};
  \node[th] at (99.5,9) {T2};
  \node[lab, anchor=south] at (105,12.5) {\texttt{not\_full}};
  \node[q] (ne) at (105,-9) {};
  \node[th] at (99.5,-9) {T3};
  \node[lab, anchor=south] at (105,-5.5) {\texttt{not\_empty}};
  \draw[->] (84,9) -- (nf.west);
  \draw[->] (84,-9) -- (ne.west);
  \node[lab, anchor=south] at (89,11.5) {\texttt{wait}};
  \node[lab] at (90,0) {\iflangja{ロック\\を解放}{releases\\lock}};
  % signal path back to the entry queue
  \draw[->] (ne.south) -- ++(0,-12) -| node[lab, pos=0.25, below]
    {\texttt{signal} / \texttt{broadcast}:
     \iflangja{起こされたスレッドはロックを再獲得する}{a woken thread re-acquires the lock}} (eq.south);
  % exit
  \draw[->] (58,19) -- ++(0,6) node[lab, anchor=south] {\iflangja{return（ロックを解放）}{return (releases lock)}};
\end{tikzpicture}

  \caption{A monitor.  Threads enter only through entry procedures, which
    acquire the monitor's lock implicitly; a thread that waits on a
    condition variable releases the lock and must re-acquire it after being
    woken.}
  \label{fig:l10-monitor}
\end{figure}

Monitors were part of the Mesa language developed at Xerox PARC
\cite{lampson1980}.\margin{In Mesa and Java a signalled thread only becomes
eligible to re-enter the monitor, so it must re-test its condition in a
\texttt{while} loop.}  Java supports them today: every object has an
internal lock, and a method declared \texttt{synchronized} is an entry
procedure whose locking and unlocking the language runtime performs.
Notifications, however, remain explicit calls to \texttt{notify} and
\texttt{notifyAll}.

The term \emph{monitor} also denotes a programming style.  Structuring
every critical section with explicit enter and exit code built from a mutex
and condition variables, as in Lesson~3, applies the discipline of a
monitor without support from the language.

\section{Other synchronization constructs}
\label{sec:l10-other}

Several further constructs aim at making synchronization policies easier to
state, or at avoiding waiting altogether.\source{P3L4, more synchronization
constructs.}
\begin{description}
  \item[Serializers] A \term{serializer}[シリアライザ] \cite{hewitt1979}
    guards a resource like a monitor, but it makes priorities easy to
    define and hides explicit signalling.  A thread waits in a queue
    together with the condition under which it may proceed; the serializer
    resumes it automatically when the condition holds and the thread is at
    the head of its queue, so neither condition variables nor signal calls
    appear in the program, and the order of the queues expresses priority.
  \item[Path expressions] A \term{path expression}[パス式]
    \cite{campbell1974} states, as a regular expression over the names of
    the operations on a resource, the orders in which the operations may
    execute; the run-time system then interleaves the operations only in
    ways the expression allows.  The policy \enquote{many reads or a single
    write} is written \texttt{path \{read\}, write end}: the braces allow
    any number of simultaneous reads, the comma selects one of the
    alternatives, and the path repeats indefinitely.
  \item[Barriers] A \term{barrier}[バリア] is the reverse of a semaphore.
    A semaphore lets up to $n$ threads proceed and blocks those that come
    after them; a barrier blocks threads until $n$ of them have arrived and
    then lets all of them proceed together (\cref{fig:l10-barrier}).  POSIX
    threads provide \texttt{pthread\_barrier\_t} with
    \texttt{pthread\_barrier\_wait}.
  \item[Rendezvous points] A \term{rendezvous point}[ランデブーポイント] is
    a point in the execution of two or more specific threads at which each
    waits until the others have reached their corresponding points; the
    threads then continue, typically after exchanging data.
\end{description}

\begin{figure}[htbp]
  \centering
  % A barrier for n = 3 threads: threads block at the barrier until the
% third one arrives, then all continue.
% Usage: % A barrier for n = 3 threads: threads block at the barrier until the
% third one arrives, then all continue.
% Usage: % A barrier for n = 3 threads: threads block at the barrier until the
% third one arrives, then all continue.
% Usage: \input{figures/l10-barrier}   (width ~116 mm)
\begin{tikzpicture}[x=1mm, y=1mm, font=\footnotesize\sffamily,
  lab/.style={font=\scriptsize\sffamily, align=center},
  run/.style={line width=1.6pt, ln-accent},
  blk/.style={line width=4pt, ln-caution!25},
  >={Stealth[length=1.8mm]}]
  \foreach \y/\t/\a in {10/T1/40, 3/T2/62, -4/T3/82} {
    \node[lab, anchor=east] at (18,\y) {\t};
    \draw[run] (20,\y) -- (\a,\y);
    \ifnum\a<82 \draw[blk] (\a,\y) -- (82,\y); \fi
    \draw[run, ->] (82,\y) -- (112,\y);
    \fill[ln-accent] (\a,\y) circle (0.9);
  }
  \draw[line width=1.2pt] (82,15) -- (82,-9);
  \node[lab, anchor=south] at (82,15)
    {\iflangja{バリア（$n=3$）}{barrier ($n=3$)}};
  \node[lab, anchor=south] at (40,12)
    {\iflangja{T1 が到着}{T1 arrives}};
  \node[lab, anchor=west] at (83,-7.5)
    {\iflangja{T3 の到着で全員が進む}{T3 arrives: all proceed}};
  % legend and time axis
  \draw[run] (20,-15) -- (28,-15);
  \node[lab, anchor=west] at (29,-15) {\iflangja{実行中}{running}};
  \draw[blk] (46,-15) -- (54,-15);
  \node[lab, anchor=west] at (55,-15) {\iflangja{バリアで待機}{blocked at the barrier}};
  \draw[->] (92,-15) -- node[lab, above] {\iflangja{時間}{time}} (116,-15);
\end{tikzpicture}
   (width ~116 mm)
\begin{tikzpicture}[x=1mm, y=1mm, font=\footnotesize\sffamily,
  lab/.style={font=\scriptsize\sffamily, align=center},
  run/.style={line width=1.6pt, ln-accent},
  blk/.style={line width=4pt, ln-caution!25},
  >={Stealth[length=1.8mm]}]
  \foreach \y/\t/\a in {10/T1/40, 3/T2/62, -4/T3/82} {
    \node[lab, anchor=east] at (18,\y) {\t};
    \draw[run] (20,\y) -- (\a,\y);
    \ifnum\a<82 \draw[blk] (\a,\y) -- (82,\y); \fi
    \draw[run, ->] (82,\y) -- (112,\y);
    \fill[ln-accent] (\a,\y) circle (0.9);
  }
  \draw[line width=1.2pt] (82,15) -- (82,-9);
  \node[lab, anchor=south] at (82,15)
    {\iflangja{バリア（$n=3$）}{barrier ($n=3$)}};
  \node[lab, anchor=south] at (40,12)
    {\iflangja{T1 が到着}{T1 arrives}};
  \node[lab, anchor=west] at (83,-7.5)
    {\iflangja{T3 の到着で全員が進む}{T3 arrives: all proceed}};
  % legend and time axis
  \draw[run] (20,-15) -- (28,-15);
  \node[lab, anchor=west] at (29,-15) {\iflangja{実行中}{running}};
  \draw[blk] (46,-15) -- (54,-15);
  \node[lab, anchor=west] at (55,-15) {\iflangja{バリアで待機}{blocked at the barrier}};
  \draw[->] (92,-15) -- node[lab, above] {\iflangja{時間}{time}} (116,-15);
\end{tikzpicture}
   (width ~116 mm)
\begin{tikzpicture}[x=1mm, y=1mm, font=\footnotesize\sffamily,
  lab/.style={font=\scriptsize\sffamily, align=center},
  run/.style={line width=1.6pt, ln-accent},
  blk/.style={line width=4pt, ln-caution!25},
  >={Stealth[length=1.8mm]}]
  \foreach \y/\t/\a in {10/T1/40, 3/T2/62, -4/T3/82} {
    \node[lab, anchor=east] at (18,\y) {\t};
    \draw[run] (20,\y) -- (\a,\y);
    \ifnum\a<82 \draw[blk] (\a,\y) -- (82,\y); \fi
    \draw[run, ->] (82,\y) -- (112,\y);
    \fill[ln-accent] (\a,\y) circle (0.9);
  }
  \draw[line width=1.2pt] (82,15) -- (82,-9);
  \node[lab, anchor=south] at (82,15)
    {\iflangja{バリア（$n=3$）}{barrier ($n=3$)}};
  \node[lab, anchor=south] at (40,12)
    {\iflangja{T1 が到着}{T1 arrives}};
  \node[lab, anchor=west] at (83,-7.5)
    {\iflangja{T3 の到着で全員が進む}{T3 arrives: all proceed}};
  % legend and time axis
  \draw[run] (20,-15) -- (28,-15);
  \node[lab, anchor=west] at (29,-15) {\iflangja{実行中}{running}};
  \draw[blk] (46,-15) -- (54,-15);
  \node[lab, anchor=west] at (55,-15) {\iflangja{バリアで待機}{blocked at the barrier}};
  \draw[->] (92,-15) -- node[lab, above] {\iflangja{時間}{time}} (116,-15);
\end{tikzpicture}

  \caption{A barrier for three threads.  T1 and T2 block when they reach
    the barrier; the arrival of T3 releases all three.}
  \label{fig:l10-barrier}
\end{figure}

\subsection{Wait-free synchronization and RCU}

All the constructs above make threads wait.  An optimistic alternative lets
threads proceed without locks, betting that concurrent operations will not
conflict, and designs the data structure so that concurrent readers remain
safe while it is being modified.

\begin{definition}[Wait-free synchronization]
  Synchronization is \term{wait-free synchronization}[ウェイトフリー同期]
  when every thread completes each of its operations in a bounded number of
  its own steps, regardless of how fast or slow the other threads are
  \cite{herlihy1991}.  No thread ever waits for another.
\end{definition}

The \term{read-copy-update}[RCU]\index{RCU|see{read-copy-update}} (RCU)
mechanism of the Linux kernel \cite{mckenney1998} applies this idea to data
that is read far more often than it is written.  Readers access the data
without taking any lock.  A writer does not modify the data in place: it
copies the current version, modifies the copy, and publishes it by
atomically replacing the pointer through which readers find the data.
Readers that started before the replacement continue to see the old,
consistent version; readers that start afterwards see the new one.  The old
version is freed only after all readers that could still hold a reference
to it have finished.  RCU readers never wait, whereas writers still
coordinate among themselves with an ordinary lock.

\Cref{tab:l10-constructs} compares the constructs of this lesson with the
mutex.  Each of them, whether it makes threads wait or not, depends on
updating a memory location atomically---setting a lock word, adjusting a
count or replacing a pointer---and this requires support from the hardware.

\begin{table}[htbp]
  \centering
  \caption{Synchronization constructs compared.}
  \label{tab:l10-constructs}
  \small
  \begin{tabular}{@{}lp{32mm}p{30mm}l@{}}
    \toprule
    Construct & Admitted at once & Waiting thread & Locking \\
    \midrule
    Mutex              & one                         & blocks            & explicit \\
    Spinlock           & one                         & spins             & explicit \\
    Counting semaphore & up to the initial count     & blocks            & explicit \\
    Reader-writer lock & any readers or one writer   & blocks or spins   & explicit \\
    Monitor            & one                         & blocks            & implicit \\
    Barrier            & $n$, once $n$ have arrived  & blocks            & explicit \\
    RCU                & any readers; one writer     & readers never wait & writers only \\
    \bottomrule
  \end{tabular}
\end{table}

\begin{keypoint}
  Spinlocks trade CPU time for context switches; semaphores admit a bounded
  number of threads; reader-writer locks package the readers-writers
  policy; monitors make locking implicit.  Serializers, path expressions,
  barriers and rendezvous points express further policies, and wait-free
  schemes such as RCU avoid waiting altogether.
\end{keypoint}

\section{Hardware support for synchronization}
\label{sec:l10-hw}

We first try to implement a spinlock with ordinary
instructions.\source{P3L4, sync building block: spinlock, need for
hardware support, atomic instructions; OSTEP ch.~28; Silberschatz et al.\
ch.~6.}
\needspace{10\baselineskip}
\begin{lstlisting}[language=C, numbers=none]
#define FREE 0
#define BUSY 1

void naive_lock(int *lock) {
  while (*lock == BUSY)
    ;                /* spin */
  *lock = BUSY;      /* separate from the test */
}
\end{lstlisting}
The test and the update are separate instructions, and between them
another thread can observe the same free lock.  \Cref{tab:l10-naive-race}
shows two threads on different CPUs that both enter the critical section.
Protecting the two steps with another lock only moves the problem into that
lock.  Mutual exclusion requires that checking the lock and setting it
happen as one indivisible step, and only the hardware can provide such a
step.

\begin{table}[htbp]
  \centering
  \caption{Two threads executing \texttt{naive\_lock} on different CPUs.
    Both read the lock before either writes it, and both enter the
    critical section.}
  \label{tab:l10-naive-race}
  \small
  \begin{tabular}{@{}clll@{}}
    \toprule
    Step & T1 on CPU 1 & T2 on CPU 2 & \texttt{*lock} \\
    \midrule
    1 & reads \texttt{FREE}, leaves loop &                          & \texttt{FREE} \\
    2 &                         & reads \texttt{FREE}, leaves loop  & \texttt{FREE} \\
    3 & writes \texttt{BUSY}, enters    &                          & \texttt{BUSY} \\
    4 &                         & writes \texttt{BUSY}, enters      & \texttt{BUSY} \\
    \bottomrule
  \end{tabular}
\end{table}

\begin{definition}[Atomic instruction]
  An \term{atomic instruction}[アトミック命令] is a processor instruction
  that performs several operations, typically reading a memory location and
  writing it, as one indivisible step.  The hardware guarantees atomicity
  (the operations take effect together), mutual exclusion (no other
  execution of an atomic instruction interleaves with it), and that of
  several concurrent executions on the same location all but one are made
  to wait.
\end{definition}

An atomic instruction is thus a critical section protected by the hardware.
Common examples are:
\begin{description}
  \item[\texttt{test\_and\_set}] The \term{test-and-set}[テストアンドセット]
    instruction returns the current value of a lock word and sets the word
    to \texttt{BUSY}.
  \item[\texttt{read\_and\_increment}] returns the value of a memory
    location and increments it.
  \item[\texttt{compare\_and\_swap}] The
    \term{compare-and-swap}[コンペアアンドスワップ] instruction compares a
    memory location with an expected value and, only if they are equal,
    writes a new value; it reports whether the write took place.
\end{description}

\needspace{14\baselineskip}
Written as C functions that the processor executes atomically, the first
and the last are:
\begin{lstlisting}[language=C, numbers=none]
/* each function executes as one atomic step */
int test_and_set(int *lock) {
  int old = *lock;
  *lock = BUSY;
  return old;
}

int compare_and_swap(int *p, int expected,
                     int new_val) {
  if (*p != expected)
    return 0;            /* value has changed */
  *p = new_val;
  return 1;
}
\end{lstlisting}

\needspace{9\baselineskip}
With \texttt{test\_and\_set} the spinlock becomes correct:
\begin{lstlisting}[language=C, numbers=none]
void spinlock_lock(int *lock) {
  while (test_and_set(lock) == BUSY)
    ;   /* spin */
}

void spinlock_unlock(int *lock) {
  *lock = FREE;
}
\end{lstlisting}
If the lock is free, \texttt{test\_and\_set} returns \texttt{FREE} and has
already marked the lock busy, so the thread leaves the loop as the owner.
If the lock is busy, it returns \texttt{BUSY} after writing \texttt{BUSY}
over \texttt{BUSY}, which changes nothing, and the thread tries again.  Of
several threads that execute \texttt{test\_and\_set} on a free lock at the
same time, exactly one obtains \texttt{FREE}.

\begin{example}
  Compare-and-swap updates shared data without any lock.  A counter is
  incremented by re-reading it until the swap succeeds:
  \begin{lstlisting}[language=C, numbers=none]
void increment(int *counter) {
  int old;
  do {
    old = *counter;
  } while (!compare_and_swap(counter, old, old + 1));
}
  \end{lstlisting}
  The swap fails exactly when another thread changed the counter between
  the read and the swap; the loop then re-reads and retries, so no
  increment is lost.
\end{example}

\caution{The retry loop never blocks, but under contention a thread may
fail repeatedly, so it is not wait-free in the sense of
\cref{sec:l10-other}.}

\subsection{Portability}

Atomic instructions are specific to the processor architecture: which ones
exist and how efficient they are differ between platforms.  On x86, for
example, \texttt{xchg} with a memory operand serves as test-and-set and
\texttt{lock cmpxchg} as compare-and-swap.  Software that implements
synchronization constructs, such as a kernel or a thread library, must
therefore be ported to every architecture it runs on.  It is also
specialized for each platform: conditional compilation selects, among
several implementations of the same construct, the one built on the most
efficient instructions the platform offers.

\section{Shared memory multiprocessors and caches}
\label{sec:l10-smp}

Synchronization constructs matter most on machines with several CPUs,
where threads really do run at the same time.\source{P3L4, shared memory
multiprocessors, cache coherence, cache coherence and atomic instructions;
Silberschatz et al.\ ch.~1.}  A shared memory multiprocessor
(Lesson~7)\index{shared memory multiprocessor} has several CPUs and memory
that all of them can access; such systems are also called symmetric
multiprocessors (SMP).  \Cref{fig:l10-smp} shows two configurations.  In a
bus-based configuration, typical of older systems, the CPUs and a single
memory component are attached to one shared bus, which carries only one
memory access at a time.  In an interconnect-based configuration, common
today, the CPUs and possibly several memory components are connected
through an interconnect on which several memory accesses can be in
progress at the same time.

\begin{figure}[htbp]
  \centering
  % Shared memory multiprocessor configurations: bus-based and
% interconnect-based, each CPU with its own cache.
% Usage: % Shared memory multiprocessor configurations: bus-based and
% interconnect-based, each CPU with its own cache.
% Usage: % Shared memory multiprocessor configurations: bus-based and
% interconnect-based, each CPU with its own cache.
% Usage: \input{figures/l10-smp}   (width ~116 mm)
\begin{tikzpicture}[x=1mm, y=1mm, font=\footnotesize\sffamily,
  cpu/.style={draw, fill=white, minimum width=13mm, minimum height=6mm, inner sep=1pt},
  cache/.style={draw, fill=ln-box, minimum width=13mm, minimum height=4.5mm,
                inner sep=1pt, font=\scriptsize\sffamily},
  mem/.style={draw, fill=ln-box, minimum height=7mm, inner sep=1pt},
  lab/.style={font=\scriptsize\sffamily, align=center}]
  % (a) bus-based
  \foreach \x/\n in {8/1, 27/2, 46/3} {
    \node[cpu] (a\n) at (\x,0) {CPU \n};
    \node[cache, anchor=north] (ac\n) at (a\n.south) {\iflangja{キャッシュ}{cache}};
    \draw (ac\n.south) -- (\x,-13);
  }
  \draw[line width=1.6pt] (0,-13) -- (54,-13);
  \node[lab, anchor=south west] at (48,-12.6) {\iflangja{バス}{bus}};
  \node[mem, minimum width=30mm] (am) at (27,-22) {\iflangja{メモリ}{memory}};
  \draw (am.north) -- (27,-13);
  \node[lab] at (27,-31) {\iflangja{(a) バス型}{(a) bus-based}};
  % (b) interconnect-based
  \foreach \x/\n in {70/1, 89/2, 108/3} {
    \node[cpu] (b\n) at (\x,0) {CPU \n};
    \node[cache, anchor=north] (bc\n) at (b\n.south) {\iflangja{キャッシュ}{cache}};
    \draw (bc\n.south) -- (\x,-10.5);
  }
  \node[draw, fill=white, minimum width=50mm, minimum height=5mm, inner sep=1pt]
    (ic) at (89,-13) {\iflangja{インターコネクト}{interconnect}};
  \node[mem, minimum width=22mm] (bm1) at (76,-22) {\iflangja{メモリ}{memory}};
  \node[mem, minimum width=22mm] (bm2) at (102,-22) {\iflangja{メモリ}{memory}};
  \draw (bm1.north) -- (76,-15.5);
  \draw (bm2.north) -- (102,-15.5);
  \node[lab] at (89,-31) {\iflangja{(b) インターコネクト型}{(b) interconnect-based}};
\end{tikzpicture}
   (width ~116 mm)
\begin{tikzpicture}[x=1mm, y=1mm, font=\footnotesize\sffamily,
  cpu/.style={draw, fill=white, minimum width=13mm, minimum height=6mm, inner sep=1pt},
  cache/.style={draw, fill=ln-box, minimum width=13mm, minimum height=4.5mm,
                inner sep=1pt, font=\scriptsize\sffamily},
  mem/.style={draw, fill=ln-box, minimum height=7mm, inner sep=1pt},
  lab/.style={font=\scriptsize\sffamily, align=center}]
  % (a) bus-based
  \foreach \x/\n in {8/1, 27/2, 46/3} {
    \node[cpu] (a\n) at (\x,0) {CPU \n};
    \node[cache, anchor=north] (ac\n) at (a\n.south) {\iflangja{キャッシュ}{cache}};
    \draw (ac\n.south) -- (\x,-13);
  }
  \draw[line width=1.6pt] (0,-13) -- (54,-13);
  \node[lab, anchor=south west] at (48,-12.6) {\iflangja{バス}{bus}};
  \node[mem, minimum width=30mm] (am) at (27,-22) {\iflangja{メモリ}{memory}};
  \draw (am.north) -- (27,-13);
  \node[lab] at (27,-31) {\iflangja{(a) バス型}{(a) bus-based}};
  % (b) interconnect-based
  \foreach \x/\n in {70/1, 89/2, 108/3} {
    \node[cpu] (b\n) at (\x,0) {CPU \n};
    \node[cache, anchor=north] (bc\n) at (b\n.south) {\iflangja{キャッシュ}{cache}};
    \draw (bc\n.south) -- (\x,-10.5);
  }
  \node[draw, fill=white, minimum width=50mm, minimum height=5mm, inner sep=1pt]
    (ic) at (89,-13) {\iflangja{インターコネクト}{interconnect}};
  \node[mem, minimum width=22mm] (bm1) at (76,-22) {\iflangja{メモリ}{memory}};
  \node[mem, minimum width=22mm] (bm2) at (102,-22) {\iflangja{メモリ}{memory}};
  \draw (bm1.north) -- (76,-15.5);
  \draw (bm2.north) -- (102,-15.5);
  \node[lab] at (89,-31) {\iflangja{(b) インターコネクト型}{(b) interconnect-based}};
\end{tikzpicture}
   (width ~116 mm)
\begin{tikzpicture}[x=1mm, y=1mm, font=\footnotesize\sffamily,
  cpu/.style={draw, fill=white, minimum width=13mm, minimum height=6mm, inner sep=1pt},
  cache/.style={draw, fill=ln-box, minimum width=13mm, minimum height=4.5mm,
                inner sep=1pt, font=\scriptsize\sffamily},
  mem/.style={draw, fill=ln-box, minimum height=7mm, inner sep=1pt},
  lab/.style={font=\scriptsize\sffamily, align=center}]
  % (a) bus-based
  \foreach \x/\n in {8/1, 27/2, 46/3} {
    \node[cpu] (a\n) at (\x,0) {CPU \n};
    \node[cache, anchor=north] (ac\n) at (a\n.south) {\iflangja{キャッシュ}{cache}};
    \draw (ac\n.south) -- (\x,-13);
  }
  \draw[line width=1.6pt] (0,-13) -- (54,-13);
  \node[lab, anchor=south west] at (48,-12.6) {\iflangja{バス}{bus}};
  \node[mem, minimum width=30mm] (am) at (27,-22) {\iflangja{メモリ}{memory}};
  \draw (am.north) -- (27,-13);
  \node[lab] at (27,-31) {\iflangja{(a) バス型}{(a) bus-based}};
  % (b) interconnect-based
  \foreach \x/\n in {70/1, 89/2, 108/3} {
    \node[cpu] (b\n) at (\x,0) {CPU \n};
    \node[cache, anchor=north] (bc\n) at (b\n.south) {\iflangja{キャッシュ}{cache}};
    \draw (bc\n.south) -- (\x,-10.5);
  }
  \node[draw, fill=white, minimum width=50mm, minimum height=5mm, inner sep=1pt]
    (ic) at (89,-13) {\iflangja{インターコネクト}{interconnect}};
  \node[mem, minimum width=22mm] (bm1) at (76,-22) {\iflangja{メモリ}{memory}};
  \node[mem, minimum width=22mm] (bm2) at (102,-22) {\iflangja{メモリ}{memory}};
  \draw (bm1.north) -- (76,-15.5);
  \draw (bm2.north) -- (102,-15.5);
  \node[lab] at (89,-31) {\iflangja{(b) インターコネクト型}{(b) interconnect-based}};
\end{tikzpicture}

  \caption{Shared memory multiprocessor configurations.  In both, every CPU
    has its own cache and all CPUs can access all memory.}
  \label{fig:l10-smp}
\end{figure}

Every CPU has its own cache.  Memory is slower to reach than the cache,
both because it is physically farther away and because the CPUs contend
for the bus or interconnect that leads to it; the caches hide this latency.

When several CPUs work on the same data, however, copies of one memory
location may reside in several caches at once.  If one CPU writes the
location, the copies held by the other CPUs become stale, and a thread that
read such a copy would see an outdated value (\cref{fig:l10-coherence}).
The property that every CPU observes the most recent value of a memory
location, even though copies of it may be cached on several CPUs, is called
\term{cache coherence}[キャッシュ一貫性].

\begin{figure}[htbp]
  \centering
  % Cache coherence: a write by one CPU makes the cached copies of the same
% memory location on the other CPUs stale; the hardware must update or
% invalidate them.
% Usage: % Cache coherence: a write by one CPU makes the cached copies of the same
% memory location on the other CPUs stale; the hardware must update or
% invalidate them.
% Usage: % Cache coherence: a write by one CPU makes the cached copies of the same
% memory location on the other CPUs stale; the hardware must update or
% invalidate them.
% Usage: \input{figures/l10-coherence}   (width ~116 mm)
\begin{tikzpicture}[x=1mm, y=1mm, font=\footnotesize\sffamily,
  cpu/.style={draw, fill=white, minimum width=30mm, minimum height=6mm, inner sep=1pt},
  cache/.style={draw, fill=ln-box, minimum width=30mm, minimum height=6mm,
                inner sep=1pt, font=\scriptsize\sffamily},
  lab/.style={font=\scriptsize\sffamily, align=center},
  >={Stealth[length=1.8mm]}]
  \node[cpu] (c1) at (17,14) {CPU 1: \texttt{test\_and\_set}};
  \node[cpu] (c2) at (58,14) {CPU 2};
  \node[cpu] (c3) at (99,14) {CPU 3};
  \node[cache, anchor=north, fill=white, thick] (k1) at (c1.south)
    {\texttt{lock = BUSY}};
  \node[cache, anchor=north] (k2) at (c2.south)
    {\texttt{lock = FREE}\ \iflangja{（古い）}{(stale)}};
  \node[cache, anchor=north] (k3) at (c3.south)
    {\texttt{lock = FREE}\ \iflangja{（古い）}{(stale)}};
  \draw[line width=1.6pt] (0,-6) -- (116,-6);
  \foreach \k in {1,2,3} { \draw (k\k.south) -- (k\k.south |- 0,-6); }
  \node[draw, fill=ln-box, minimum width=40mm, minimum height=6mm] (m) at (58,-15)
    {\iflangja{メモリ}{memory}: \texttt{lock = BUSY}};
  \draw (m.north) -- (58,-6);
  % write path
  \draw[->, thick, ln-accent] ([xshift=-6mm]k1.south) -- ++(0,-12) |- (m.west)
    node[lab, pos=0.25, left, text=ln-accent] {\iflangja{(1) 書き込み}{(1) write}};
  % coherence
  \draw[->, thick, ln-caution, densely dashed] (m.east) -| ([xshift=6mm]k3.south);
  \node[lab, text=ln-caution] at (89,-9.5)
    {\iflangja{(2) 更新または無効化}{(2) update or invalidate}};
  \draw[->, thick, ln-caution, densely dashed] ([xshift=8mm]m.north) -- ([xshift=8mm]k2.south);
\end{tikzpicture}
   (width ~116 mm)
\begin{tikzpicture}[x=1mm, y=1mm, font=\footnotesize\sffamily,
  cpu/.style={draw, fill=white, minimum width=30mm, minimum height=6mm, inner sep=1pt},
  cache/.style={draw, fill=ln-box, minimum width=30mm, minimum height=6mm,
                inner sep=1pt, font=\scriptsize\sffamily},
  lab/.style={font=\scriptsize\sffamily, align=center},
  >={Stealth[length=1.8mm]}]
  \node[cpu] (c1) at (17,14) {CPU 1: \texttt{test\_and\_set}};
  \node[cpu] (c2) at (58,14) {CPU 2};
  \node[cpu] (c3) at (99,14) {CPU 3};
  \node[cache, anchor=north, fill=white, thick] (k1) at (c1.south)
    {\texttt{lock = BUSY}};
  \node[cache, anchor=north] (k2) at (c2.south)
    {\texttt{lock = FREE}\ \iflangja{（古い）}{(stale)}};
  \node[cache, anchor=north] (k3) at (c3.south)
    {\texttt{lock = FREE}\ \iflangja{（古い）}{(stale)}};
  \draw[line width=1.6pt] (0,-6) -- (116,-6);
  \foreach \k in {1,2,3} { \draw (k\k.south) -- (k\k.south |- 0,-6); }
  \node[draw, fill=ln-box, minimum width=40mm, minimum height=6mm] (m) at (58,-15)
    {\iflangja{メモリ}{memory}: \texttt{lock = BUSY}};
  \draw (m.north) -- (58,-6);
  % write path
  \draw[->, thick, ln-accent] ([xshift=-6mm]k1.south) -- ++(0,-12) |- (m.west)
    node[lab, pos=0.25, left, text=ln-accent] {\iflangja{(1) 書き込み}{(1) write}};
  % coherence
  \draw[->, thick, ln-caution, densely dashed] (m.east) -| ([xshift=6mm]k3.south);
  \node[lab, text=ln-caution] at (89,-9.5)
    {\iflangja{(2) 更新または無効化}{(2) update or invalidate}};
  \draw[->, thick, ln-caution, densely dashed] ([xshift=8mm]m.north) -- ([xshift=8mm]k2.south);
\end{tikzpicture}
   (width ~116 mm)
\begin{tikzpicture}[x=1mm, y=1mm, font=\footnotesize\sffamily,
  cpu/.style={draw, fill=white, minimum width=30mm, minimum height=6mm, inner sep=1pt},
  cache/.style={draw, fill=ln-box, minimum width=30mm, minimum height=6mm,
                inner sep=1pt, font=\scriptsize\sffamily},
  lab/.style={font=\scriptsize\sffamily, align=center},
  >={Stealth[length=1.8mm]}]
  \node[cpu] (c1) at (17,14) {CPU 1: \texttt{test\_and\_set}};
  \node[cpu] (c2) at (58,14) {CPU 2};
  \node[cpu] (c3) at (99,14) {CPU 3};
  \node[cache, anchor=north, fill=white, thick] (k1) at (c1.south)
    {\texttt{lock = BUSY}};
  \node[cache, anchor=north] (k2) at (c2.south)
    {\texttt{lock = FREE}\ \iflangja{（古い）}{(stale)}};
  \node[cache, anchor=north] (k3) at (c3.south)
    {\texttt{lock = FREE}\ \iflangja{（古い）}{(stale)}};
  \draw[line width=1.6pt] (0,-6) -- (116,-6);
  \foreach \k in {1,2,3} { \draw (k\k.south) -- (k\k.south |- 0,-6); }
  \node[draw, fill=ln-box, minimum width=40mm, minimum height=6mm] (m) at (58,-15)
    {\iflangja{メモリ}{memory}: \texttt{lock = BUSY}};
  \draw (m.north) -- (58,-6);
  % write path
  \draw[->, thick, ln-accent] ([xshift=-6mm]k1.south) -- ++(0,-12) |- (m.west)
    node[lab, pos=0.25, left, text=ln-accent] {\iflangja{(1) 書き込み}{(1) write}};
  % coherence
  \draw[->, thick, ln-caution, densely dashed] (m.east) -| ([xshift=6mm]k3.south);
  \node[lab, text=ln-caution] at (89,-9.5)
    {\iflangja{(2) 更新または無効化}{(2) update or invalidate}};
  \draw[->, thick, ln-caution, densely dashed] ([xshift=8mm]m.north) -- ([xshift=8mm]k2.south);
\end{tikzpicture}

  \caption{CPU 1 sets \texttt{lock} to \texttt{BUSY}.  The copies in the
    caches of CPU 2 and CPU 3 become stale; on a cache-coherent platform
    the hardware updates or invalidates them.}
  \label{fig:l10-coherence}
\end{figure}

Platforms differ in who is responsible for cache coherence.  On a
non-cache-coherent (NCC) platform the hardware does not propagate a write
to the copies in other caches, and software must deal with stale copies.
On a cache-coherent (CC) platform the hardware keeps the caches coherent:
it updates the other copies or invalidates them so that the next access
fetches the value from memory.  Even then the cost of the coherence
mechanism is incurred by software that writes shared locations.

\subsection{Atomic instructions on multiprocessors}

On a multiprocessor an atomic instruction must be atomic across all CPUs:
its execution on one CPU must exclude executions on the same location on
every other CPU.  It therefore cannot work on a private cached copy; it
must act on the shared memory location and leave every cache coherent.
Atomic instructions consequently take longer than ordinary instructions on
the same data: they must reach memory, contending with the other CPUs for
the bus or interconnect, and they set off the coherence mechanism.

Spinlocks make this worse.  \texttt{test\_and\_set} always writes, even
when it only replaces \texttt{BUSY} with \texttt{BUSY}, and every write
starts the cache coherence mechanism.  A spinning thread executes
\texttt{test\_and\_set} continuously, so a spinlock that several threads
wait for generates a steady stream of coherence updates on every CPU that
holds a copy of the lock word.

\begin{keypoint}
  Every synchronization construct rests on atomic instructions such as
  test-and-set and compare-and-swap.  On a shared memory multiprocessor
  they must be atomic across all CPUs and keep all caches coherent, which
  makes them, and especially repeated writes to a contended lock word,
  expensive.
\end{keypoint}

\section{Spinlock performance}
\label{sec:l10-spin-perf}

The cost of atomic instructions on multiprocessors affects how a spinlock
should be written, a question studied by Anderson et
al.\ \cite{anderson1989}.\source{P3L4, spinlock performance comparisons;
Anderson et al.\ \cite{anderson1989}.}  Three metrics characterize a
spinlock implementation:
\begin{description}
  \item[Latency] the time a thread needs to acquire a lock that is free.
    Ideally the thread executes the atomic instruction immediately.
  \item[Delay] the time between the release of the lock and its
    acquisition by a spinning thread.  Ideally a spinning thread notices
    the release immediately, which requires it to test the lock
    continuously.
  \item[Contention] the overhead caused by
    \term{lock contention}[ロック競合], the situation in which several
    threads try to acquire the same lock at the same time; above all, the
    cache coherence traffic caused by their atomic instructions.  Ideally a
    spinning thread writes nothing while the lock is held.
\end{description}

The \texttt{test\_and\_set} loop of \cref{sec:l10-hw} is ideal in latency
and delay but poor in contention, because it writes on every iteration.
The
\term{test-and-test-and-set}[テストアンドテストアンドセット] technique
first reads the lock with an ordinary instruction and executes
\texttt{test\_and\_set} only when the lock appears free:
\needspace{5\baselineskip}
\begin{lstlisting}[language=C, numbers=none]
void spinlock_lock(int *lock) {
  while (*lock == BUSY || test_and_set(lock) == BUSY)
    ;   /* spin */
}
\end{lstlisting}
The short-circuit \texttt{||} evaluates \texttt{test\_and\_set} only when
the read found \texttt{FREE}.  The read alone does not decide
anything---another thread may acquire the lock between the read and the
\texttt{test\_and\_set}---but in that case the atomic
\texttt{test\_and\_set} returns \texttt{BUSY} and the loop continues, so
the lock remains correct.  While the lock is held, a spinning thread
performs only reads, which cause no coherence traffic; writes occur only
around the moment of a release.  The price is one extra read when the lock
is free, a slight increase in latency, while delay is unchanged because
the loop still tests continuously.

\begin{example}
  On four CPUs, a thread on CPU~1 holds a lock for \SI{1}{\milli\second}
  while threads on CPUs 2--4 spin, each iteration of a spin loop taking
  \SI{50}{\nano\second}.  Each spinning thread performs
  $\SI{1}{\milli\second}/\SI{50}{\nano\second} = \num{20000}$ iterations.
  With the plain \texttt{test\_and\_set} loop these are $3\times\num{20000}
  = \num{60000}$ writes to the lock word, each of which forces the other
  caches to be updated or invalidated.  With test-and-test-and-set the
  \num{60000} iterations are reads; when the lock is released each of the
  three threads issues about one \texttt{test\_and\_set}, one of which
  succeeds, so only about three writes result.
\end{example}

\begin{keypoint}[Lesson summary]
  Mutexes and condition variables are error-prone, limited in expressive
  power and dependent on hardware support.  Spinlocks busy-wait instead of
  blocking and pay off only for short critical sections on multiple CPUs;
  semaphores admit up to a given number of threads, and a binary semaphore
  acts as a mutex; reader-writer locks provide shared and exclusive modes;
  monitors acquire and release their lock implicitly in entry procedures.
  Serializers, path expressions, barriers, rendezvous points and wait-free
  schemes such as RCU complete the picture.  All of them depend on atomic
  instructions such as test-and-set and compare-and-swap, which on a shared
  memory multiprocessor must also keep the CPUs' caches coherent; reading a
  lock before attempting test-and-set avoids most of the resulting write
  traffic under contention.
\end{keypoint}

\end{document}
